% POLYDIM_DEST
% destino: polydim/docs/
% nombre:  polydim_paper_arxiv_final.tex
% autor:   ai.mpat.agt@gmail.com
% fecha:   2026-06-28
% estado:  SUBMISSION READY — compilar, revisar PDF, subir a arXiv
% base:    polydim_paper_arxiv.tex + H2 Remark + 2 nuevas citas
% compilar: pdflatex polydim_paper_arxiv_final.tex
%           bibtex   polydim_paper_arxiv_final
%           pdflatex polydim_paper_arxiv_final.tex
%           pdflatex polydim_paper_arxiv_final.tex  (2x para referencias cruzadas)
% figuras:  polydim_figures.tex debe compilarse por separado para generar PDFs:
%           pdflatex polydim_figures.tex → genera 4 figuras en páginas separadas
%           Convertir cada página a PDF y subir como figure1.pdf ... figure4.pdf

\documentclass[10pt,twocolumn]{article}

%% ============================================================
%% Packages
%% ============================================================
\usepackage[utf8]{inputenc}
\usepackage[T1]{fontenc}
\usepackage[margin=1in]{geometry}
\usepackage{amsmath,amssymb,amsthm}
\usepackage{mathtools}
\usepackage{bm}
\usepackage{microtype}
\usepackage{hyperref}
\usepackage{cleveref}
\usepackage{booktabs}
\usepackage{listings}
\usepackage{xcolor}
\usepackage{graphicx}
\usepackage{enumitem}
\usepackage{authblk}
\usepackage{natbib}
\usepackage{stmaryrd}   % for \llbracket \rrbracket

%% ============================================================
%% Colors
%% ============================================================
\definecolor{polyblue}{RGB}{30,80,160}
\definecolor{codegray}{RGB}{248,248,248}
\definecolor{codegreen}{RGB}{0,128,0}
\definecolor{codered}{RGB}{160,0,0}
\definecolor{codepurple}{RGB}{120,0,120}

%% ============================================================
%% Listings
%% ============================================================
\lstdefinelanguage{polydim}{
  keywords={COMPOSE,MIX,FIXPOINT,COMPILE,RENDER,EXPORT,ATTEND,RECUR,BIND,
            position,transform,output,bind},
  keywordstyle=\color{polyblue}\bfseries,
  ndkeywords={DIM_FLUTTER,DIM_RUST,DIM_WASM,DIM_SQL,DIM_PYTHON,
              DIM_GRAPH,DIM_VECTOR,DIM_TIME,DIM_ERROR,DIM_META},
  ndkeywordstyle=\color{codepurple},
  sensitive=true,
  comment=[l]{--},
  commentstyle=\color{codegreen}\itshape,
  morestring=[b]",
}

\lstset{
  basicstyle=\ttfamily\small,
  backgroundcolor=\color{codegray},
  frame=single, framesep=3pt,
  numbers=none, breaklines=true, captionpos=b,
  tabsize=2, showstringspaces=false,
}

%% ============================================================
%% Theorems
%% ============================================================
\theoremstyle{plain}
\newtheorem{theorem}{Theorem}[section]
\newtheorem{corollary}[theorem]{Corollary}
\newtheorem{lemma}[theorem]{Lemma}
\newtheorem{proposition}[theorem]{Proposition}

\theoremstyle{definition}
\newtheorem{definition}[theorem]{Definition}
\newtheorem{example}[theorem]{Example}

\theoremstyle{remark}
\newtheorem{remark}[theorem]{Remark}

%% ============================================================
%% Math macros
%% ============================================================
\newcommand{\RR}{\mathbb{R}}
\newcommand{\Compose}{\mathsf{COMPOSE}}
\newcommand{\Mix}{\mathsf{MIX}}
\newcommand{\Fixpoint}{\mathsf{FIXPOINT}}
\newcommand{\Project}{\mathsf{PROJECT}}
\newcommand{\Attend}{\mathsf{ATTEND}}
\newcommand{\Recur}{\mathsf{RECUR}}
\newcommand{\Align}{\mathsf{ALIGN}}
\newcommand{\DimX}{\mathsf{DIM}_X}
\newcommand{\ProxX}{\Pi_X}
\newcommand{\EmbX}{\mathrm{embed}_X}
\newcommand{\ProjX}{\mathrm{proj}_X}
\newcommand{\GeoID}{\mathsf{GEO\_ID}}
\newcommand{\norm}[1]{\left\|#1\right\|}
\newcommand{\lfp}{\mathrm{lfp}}

%% ============================================================
%% Title
%% ============================================================
\title{\textbf{POLYDIM: A Transformer-Native Algebraic Programming Language\\
for High-Dimensional Latent Computing}}

\author[1]{POLYDIM Research Group}
\affil[1]{\texttt{ai.mpat.agt@gmail.com}}
\date{2026}

%% ============================================================
\begin{document}
%% ============================================================

\maketitle

%% ============================================================
\begin{abstract}
Modern artificial intelligence architectures---Transformers and State Space
Models---perform computation through high-dimensional geometric transformations
in latent spaces of thousands to tens of thousands of dimensions. Yet these
architectures are programmed and forced to communicate via discrete sequential
symbol streams, collapsing the rich geometry of their latent states to a
one-dimensional text channel. We define this as the \emph{Impedance Mismatch}:
a paradigm gap requiring constant serialization-deserialization cycles,
producing semantic loss with every AI-to-AI interaction.

We propose \textsc{Polydim}, a programming language whose fundamental unit of
computation is a geometric transformation $T\colon \RR^N \to \RR^N$. A closed
algebraic core of four operators---$\Compose$, $\Mix$, $\Fixpoint$, and
$\Project$---enables native latent-space execution, zero-serialization AI-to-AI
communication via cross-attention semantics, and a dimensional type system where
types emerge from orthogonal projections rather than declarations. We establish
four algebraic theorems: associativity of composition, linearity of
superposition, \emph{functoriality of projection} (Theorem~3, the central
contribution), and uniqueness of fixed points. The proof of Theorem~3 is
grounded in the Invariant Subspace Necessity theorem of~\citet{suarez2026invariant}
and the categorical framework of~\citet{oneill2025endofunctor}, showing that
the subspace-preservation property required by the functor proof is a direct
consequence of the transformer architecture itself. The specification
incorporates Low-Rank Adaptation (LoRA) serialization and Homotopy Type
Theory~(HoTT) for model-agnostic identity. A Python bootstrap achieves $29/29$
tests and a $64.5\%$ semantic improvement over deterministic baselines.
\end{abstract}

%% ============================================================
\section{Introduction}
%% ============================================================

\subsection{The Impedance Mismatch Problem}

Consider a large language model with representation space $\RR^{12288}$.
When two such models communicate via the standard protocol, their geometric
states must be collapsed to token sequences, transmitted, and re-expanded ---
two irreversible projections, two semantic losses, no fidelity guarantees.

\paragraph{The Gesture vs.\ Emoji Analogy.}
An AI operating natively works with a ``gesture''---a continuous manifold of
meaning in $1{,}000$+ dimensions. Text communication collapses it to a
flat ``emoji'': recognizable, inherently lossy. \textsc{Polydim} allows AIs
to exchange the full gesture as tensor transformations.

\begin{quote}\small
\textit{1950s: CPUs already performed binary arithmetic.
$\to$ Assembly formalized how to program them deliberately.}\\[4pt]
\textit{2020s: Transformers already operate in high-dimensional geometry.
$\to$ \textsc{Polydim} formalizes how to program them deliberately.}
\end{quote}

\subsection{Contributions}

\begin{enumerate}[leftmargin=*,label=\arabic*.,noitemsep]
\item \textbf{Language specification} with BNF grammar, Big-Step operational
      semantics, and denotational semantics.
\item \textbf{Two-layer architecture}: invariant algebraic core separating
      architecture-agnostic operations from transformer/SSM-specific primitives.
\item \textbf{Dimensional type system}: types emerge from projection,
      formalized as a functor $F\colon \mathcal{G} \to \mathcal{E}$
      (Theorem~\ref{thm:functor}).
\item \textbf{Zero-serialization AI$\leftrightarrow$AI protocol}: $\Align$
      via Procrustes.
\item \textbf{HoTT-based identity}: $\GeoID$ as a homotopy equivalence class.
\item \textbf{Literature grounding}: H2 derived from~\citet{suarez2026invariant}
      and~\citet{oneill2025endofunctor}.
\item \textbf{Python bootstrap}: $29/29$ tests, $0.9993$ align score,
      $+64.5\%$ semantic improvement.
\end{enumerate}

%% ============================================================
\section{Background}
%% ============================================================

\subsection{Transformers as Geometric Computation}

The central operation of a Transformer~\citep{vaswani2017}:
\[
  \mathrm{Attention}(Q,K,V) = \mathrm{softmax}\!\left(\tfrac{QK^\top}{\sqrt{d}}\right)V
\]
is $T\colon \RR^d \to \RR^d$. Every layer is
$s_{\mathrm{out}} = T_L \circ \cdots \circ T_1(s_{\mathrm{in}})$, i.e.,
$\Compose$ applied $L$ times.

\subsection{State Space Models}

Mamba~\citep{gu2023mamba}: $h(t{+}1) = Ah(t){+}Bx(t)$, $y(t){=}Ch(t)$.
A \texttt{.polydim} file can literally be the serialized matrices $(A,B,C)$.

\subsection{VSA and Quasi-Orthogonality}

For large $N$, random vectors in $\RR^N$ satisfy
$\mathbb{E}[\cos(u,v)]=0$, $\sigma[\cos(u,v)]=1/\sqrt{N}$.
At $N=10{,}000$: $\sigma\approx0.01$, so two branches of $\Mix$ coexist
without destructive interference~\citep{kanerva2009,frady2021}.

\subsection{Invariant Subspaces in Transformers}

\citet{suarez2026invariant} prove that any semantic feature decoded through a
linear interface occupies a context-invariant subspace (Theorem~3.7, ibid.).
The OV circuit maps feature direction $d_f \mapsto W_O W_V d_f$, preserving
the span of feature directions (Corollary~B.3, ibid.).

\subsection{Self-Attention as a Categorical Endofunctor}

\citet{oneill2025endofunctor} establish that self-attention induces a
parametric endofunctor $F\colon \mathrm{Vect} \to \mathrm{Vect}$ in the
2-category $\mathrm{Para}(\mathrm{Vect})$ (Theorem~3.1, ibid.), and that
stacking $n$ attention layers corresponds to the free monad $\mathrm{Free}(F)$
(Theorem~3.2, ibid.). This categorical structure is precisely what
$\Compose$ and $\Project$ formalize in \textsc{Polydim}.

%% ============================================================
\section{POLYDIM Language Design}
%% ============================================================

\subsection{Core Principles}

\textbf{P1} No instructions, only transformations ($T\colon\RR^N\to\RR^N$).
\textbf{P2} No variables, only positions.
\textbf{P3} No loops, only convergence ($\Fixpoint$).
\textbf{P4} No binary branching, only superposition ($\Mix$).

\subsection{Formal Definitions and Grammar}

\begin{definition}[State]
$S=(V,D,A)$: $V\in\RR^N$; $D\subseteq\{\mathsf{DIM}_0,\ldots,\mathsf{DIM}_8\}$;
$A\colon D\to[0,1]$.
\end{definition}

\begin{definition}[Position]
$P=(g,S)$: $g\in\RR^N$ is the $\GeoID$ (invariant); $S$ is mutable.
\end{definition}

\begin{definition}[Admissible Transformation]
$T\colon\RR^N\to\RR^N$ measurable, $\GeoID$-preserving,
and representable as $T=W_0+UV^\top$ with $U,V\in\RR^{N\times r}$, $r\ll N$.
\end{definition}

\paragraph{BNF Grammar.}
\begin{lstlisting}[language=polydim,caption={Abstract BNF.}]
<program>   ::= <transform>
<transform> ::=
    COMPOSE ( <transform> , <transform> )
  | MIX ( <w> , <transform> , <w> , <transform> )
  | FIXPOINT ( <transform> , <eps> )
  | COMPILE ( <transform> , <executor> )
  | RENDER  ( <transform> , DIM_FLUTTER )
  | EXPORT  ( <transform> , <executor> )
  | ATTEND ( <mat> , <mat> , <mat> )
  | RECUR  ( <mat> , <mat> , <mat> )
  | <position-ref>
<executor> ::= DIM_SQL | DIM_RUST | DIM_PYTHON
             | DIM_GRAPH | DIM_WASM
             | <executor> ∩ <executor>
\end{lstlisting}
The grammar excludes named variables (P2), loops (P3), binary branching (P4),
and declared types---defining absences, not restrictions.

\paragraph{Concrete Syntax.}
\begin{lstlisting}[language=polydim]
position query : R^10000
transform attend := ATTEND(W_Q, W_K, W_V)
transform refine := MIX(0.8, attend, 0.2,
                        ATTEND(W_Q2,W_K2,W_V2))
transform stable := FIXPOINT(refine, 1e-4)
output ui   := RENDER(stable, DIM_FLUTTER)
output db   := EXPORT(stable, DIM_SQL)
output live := RENDER(stable, DIM_FLUTTER ∩ DIM_SQL)
\end{lstlisting}

\subsection{Two-Layer Architecture}

\textbf{Algebraic Layer} (invariant across architectures):
\begin{align*}
\Compose(T_1,T_2)          &= T_2\circ T_1 \\
\Mix(\alpha,T_1,\beta,T_2) &= \alpha T_1+\beta T_2 \\
\Fixpoint(T,\varepsilon)   &= T^n(v)\ \text{until}\ \norm{s_{k+1}-s_k}<\varepsilon \\
\Project(T,E)              &= F_E(T(s))
\end{align*}

\textbf{Implementation Layer} (architecture-specific):
\begin{align*}
\Attend(W_Q,W_K,W_V,s) &= \mathrm{softmax}\!\left(\tfrac{(sW_Q)(sW_K)^\top}{\sqrt{d}}\right)\!(sW_V) \\
\Recur(A,B,C,h,x)      &= Ah{+}Bx,\quad y=Ch
\end{align*}

\subsection{Big-Step Operational Semantics}

\begin{gather*}
\frac{\langle s,T_1\rangle\Downarrow s_1\quad\langle s_1,T_2\rangle\Downarrow s_2}
     {\langle s,\Compose(T_1,T_2)\rangle\Downarrow s_2}\tag{E-Compose}\\[4pt]
\frac{}{\langle s,\Mix(\alpha,T_1,\beta,T_2)\rangle\Downarrow
        \alpha T_1(s){+}\beta T_2(s)}\tag{E-Mix}\\[4pt]
\frac{\exists n:\norm{T^n(s)-T^{n-1}(s)}<\varepsilon,\;s^*=T^n(s)}
     {\langle s,\Fixpoint(T,\varepsilon)\rangle\Downarrow s^*}\tag{E-Fix}\\[4pt]
\frac{F_E:\mathcal{G}\to\mathcal{E}}
     {\langle s,\Project(T,E)\rangle\Downarrow F_E(T(s))}\tag{E-Proj}
\end{gather*}

\subsection{Formal Theorems}
\label{sec:theorems}

\begin{theorem}[Associativity of $\Compose$]
\label{thm:assoc}
$\Compose(\Compose(T_3,T_2),T_1) = \Compose(T_3,\Compose(T_2,T_1))$.
\end{theorem}
\begin{proof}
Standard associativity of function composition. \qed
\end{proof}

\begin{remark}
Non-commutativity ($\Compose(T_2,T_1)\neq\Compose(T_1,T_2)$ in general)
encodes causal order topologically in the algebra.
\end{remark}

\begin{theorem}[Linearity of $\Mix$]
\label{thm:linear}
If $T_1,T_2\in\mathrm{Lin}(\RR^N)$, then $\Mix(\alpha,T_1,\beta,T_2)
\in\mathrm{Lin}(\RR^N)$.
\end{theorem}
\begin{proof}
A linear combination of linear maps is linear. \qed
\end{proof}

\begin{remark}[VSA Foundation]
Quasi-orthogonality ($\sigma\approx0.01$ for $N=10{,}000$) ensures two
simultaneous branches coexist without destructive interference, making $\Mix$
a sound replacement for binary branching~(P4).
\end{remark}

\begin{theorem}[$\Project$ is a Functor]
\label{thm:functor}
Let $\mathcal{G}=(\mathrm{Pos},\mathrm{Adm})$ be the \emph{geometric category}.
Let $\mathcal{E}_X$ be the \emph{executor category} of $\DimX$.

\begin{description}[leftmargin=1.5em,noitemsep]
\item[\textbf{H1}] $\EmbX\colon\mathrm{Ob}(\mathcal{E}_X)\to\RR^N$ satisfies
  $\ProjX(\EmbX(t))=t$ for all $t$, where
  $\ProjX(v)=\arg\min_t\norm{v-\EmbX(t)}_2$.
\item[\textbf{H2}] $T\in\mathrm{Adm}$ is $\DimX$-\emph{preserving} if
  $\ProxX T\ProxX=\ProxX T$. This holds for all $\Attend$ operators trained
  on $\DimX$-specific data by the Invariant Subspace Necessity theorem~\textnormal{\citep{suarez2026invariant}},
  and for LoRA transforms by linearity of the LoRA operator.
\end{description}

Let $T_X\subseteq\mathrm{Adm}$ denote $\DimX$-preserving transformations.
Define $\Project_X\colon\mathcal{G}|_{T_X}\to\mathcal{E}_X$ by:
\[
  \Project_X(P)=\ProjX(\ProxX V_P),\quad
  \Project_X(T)=\ProjX\circ(\ProxX T|_{\DimX})\circ\EmbX.
\]
Under H1 and H2, $\Project_X$ is a functor satisfying:\\
\textnormal{(F1)} $\Project_X(\mathrm{id}_\mathcal{G})=\mathrm{id}_{\mathcal{E}_X}$\\
\textnormal{(F2)} $\Project_X(T_2\circ T_1)=\Project_X(T_2)\circ\Project_X(T_1)$
  for all $T_1,T_2\in T_X$.
\end{theorem}

\begin{proof}
\textbf{(F1)} $\EmbX(t)\in\DimX$ by design, so $\ProxX(\EmbX(t))=\EmbX(t)$.
Then $\Project_X(\mathrm{id})(t)=\ProjX(\EmbX(t))=t$ by~H1.

\textbf{(F2)} Let $T_1,T_2\in T_X$, $t\in\mathrm{Ob}(\mathcal{E}_X)$,
$v=\EmbX(t)\in\DimX$.
By H2 for $T_1$: $T_1(v)=w_1\in\DimX$. By H2 for $T_2$:
$T_2(w_1)=w_2\in\DimX$. Since $w_2\in\DimX$, $\ProxX(w_2)=w_2$.
Therefore $\Project_X(T_2\circ T_1)(t)=\ProjX(T_2(T_1(v)))$.
Setting $t_1=\Project_X(T_1)(t)$: H1 gives $\EmbX(t_1)=T_1(v)$,
so $\Project_X(T_2)(t_1)=\ProjX(T_2(T_1(v)))$. \qed
\end{proof}

%% ── THE KEY REMARK (H2 justified from literature) ─────────
\begin{remark}[H2 Derived from Architecture]
\label{rem:h2}
Hypothesis~H2 ($\DimX$-preserving transformations) is not an ad hoc
assumption. The Invariant Subspace Necessity theorem of
\citet{suarez2026invariant} establishes that any communicable semantic feature
decoded through a linear interface must occupy a context-invariant linear
subspace (Theorem~3.7, ibid.). When $\DimX$ is defined as the span of such
invariant directions, the OV circuit maps feature directions as
$d_f \mapsto W_O W_V d_f$ (Corollary~B.3, ibid.), which remains in the same
subspace for networks trained on $\DimX$-specific data, yielding H2 for all
$\Attend$ operators. For LoRA transforms $T=W_0+UV^\top$, the same conclusion
follows by linearity of the LoRA operator over the trained subspace.
Furthermore, \citet{oneill2025endofunctor} establish that self-attention
induces an endofunctor $F\colon\mathrm{Vect}\to\mathrm{Vect}$ in
$\mathrm{Para}(\mathrm{Vect})$ (Theorem~3.1, ibid.) and that stacking attention
layers corresponds to the free monad $\mathrm{Free}(F)$ (Theorem~3.2, ibid.).
This categorical structure is precisely what $\Compose$ and $\Project_X$
formalize: the composition of transformer layers \emph{is} the free monad;
its projection to any executor is the functor $\Project_X$.
Theorem~\ref{thm:functor} is therefore grounded in the architecture of
transformers, not in a geometric axiom.
\end{remark}
%%─────────────────────────────────────────────────────────────

\begin{corollary}[Cross-Executor Portability]
\label{cor:portability}
For $\Pi=T_n\circ\cdots\circ T_1$ with all $T_i\in T_X$:
$\Project_X(\Pi)=\Project_X(T_n)\circ\cdots\circ\Project_X(T_1)$.
\end{corollary}

\begin{corollary}[Intersection as Pullback]
\label{cor:pullback}
$\Project(P,\mathsf{DIM}_A\cap\mathsf{DIM}_B)$ is the pullback of
$\Project(P,\mathsf{DIM}_A)$ and $\Project(P,\mathsf{DIM}_B)$ when
$\mathsf{DIM}_A\cap\mathsf{DIM}_B\neq\emptyset$.
This produces a live-bound type (e.g., a Flutter widget synchronized with
its SQL source) with no analogue in single-paradigm languages.
\end{corollary}

\begin{theorem}[Uniqueness of $\Fixpoint$]
\label{thm:fixpoint}
If $T$ is a contraction ($k<1$), then $\exists!\,s^*$ with $T(s^*)=s^*$
and $\Fixpoint(T,\varepsilon)$ converges to $s^*$ from any $s_0$.
\end{theorem}
\begin{proof}
Immediate from the Banach Fixed Point Theorem on $(\RR^N,\norm{\cdot}_2)$. \qed
\end{proof}

\paragraph{Denotational Semantics.}
$\llbracket\cdot\rrbracket:\mathit{Syntax}\to(\mathit{Env}\to\mathit{Val})$:
\begin{align*}
\llbracket\Compose(T_1,T_2)\rrbracket\rho
  &=\llbracket T_2\rrbracket\rho\circ\llbracket T_1\rrbracket\rho\\
\llbracket\Mix(\alpha,T_1,\beta,T_2)\rrbracket\rho
  &=\alpha\llbracket T_1\rrbracket\rho+\beta\llbracket T_2\rrbracket\rho\\
\llbracket\Fixpoint(T,\varepsilon)\rrbracket\rho
  &=\lfp_\varepsilon(\llbracket T\rrbracket\rho)\\
\llbracket\Project(T,E)\rrbracket\rho
  &=F_E\circ\llbracket T\rrbracket\rho
\end{align*}
Functoriality of $\llbracket\mathsf{COMPILE}\rrbracket$ follows from
Theorem~\ref{thm:functor}.

\subsection{Illustrative Programs}

\begin{example}[Semantic Search with Live Interface]
\begin{lstlisting}[language=polydim]
position query : R^10000
transform attend := ATTEND(W_Q, W_K, W_V)
transform refine :=
    MIX(0.8, attend, 0.2, ATTEND(W2, K2, V2))
transform stable := FIXPOINT(refine, 1e-4)
output ui   := RENDER(stable, DIM_FLUTTER)
output db   := EXPORT(stable, DIM_SQL)
-- Pullback type (Corollary 2):
output live := RENDER(stable,
                 DIM_FLUTTER ∩ DIM_SQL)
\end{lstlisting}
\texttt{live} has a type that does not exist in any single-paradigm language:
a Flutter widget synchronized live with its SQL source (Corollary~\ref{cor:pullback}).
\end{example}

\begin{example}[Zero-Serialization AI-to-AI Delegation]
\begin{lstlisting}[language=polydim]
transform task :=
    COMPOSE(
        ATTEND(W_Q, W_K, W_V),
        MIX(0.9, COMPILE(ATTEND(W2,K2,V2),
                         DIM_RUST),
            0.1, RECUR(A, B, C)))
-- ALIGN: T_B = M * task * M†  (Procrustes)
-- IA_B: new_state = T_B(state_B)
output delegation := COMPILE(task, DIM_RUST)
\end{lstlisting}
Versus text (MCP/A2A): $\RR^{12288}\to100\text{ tokens}\to\RR^{8192}$
(2 lossy projections). \textsc{Polydim} ALIGN: $T_B=MT_AM^\dagger$ (0 projections).
\end{example}

\subsection{Dimensional Type System}

A type in \textsc{Polydim} is not declared --- it emerges:
\begin{align*}
\Project(P,\mathsf{DIM\_SQL})     &\to \texttt{Column(Integer,\,NOT NULL)}\\
\Project(P,\mathsf{DIM\_FLUTTER}) &\to \texttt{TextField(value:\,42)}\\
\Project(P,\mathsf{DIM\_VECTOR})  &\to \mathbf{v}=[1.0,2.0,3.0]\in\RR^3\\
\Project(P,\mathsf{DIM\_SQL}\cap\mathsf{DIM\_FLUTTER})
  &\to \texttt{LiveDataTextField}
\end{align*}
The intersection type is the categorical pullback of Corollary~\ref{cor:pullback}.

%% ============================================================
\section{AI-to-AI Communication}
%% ============================================================

\subsection{Cross-Attention as Native Communication}

\[
  \mathrm{Communication}(s_B,s_A)
  = \mathrm{softmax}\!\left(\tfrac{Q_B K_A^\top}{\sqrt{d}}\right)V_A
\]
Reading a message from $\mathrm{IA}_A$ is mathematically identical to
computing a cross-attention layer --- not a metaphor but the exact operation
occurring in multi-encoder transformers~\citep{oneill2025endofunctor}.

\subsection{The ALIGN Protocol}

For $\mathrm{IA}_A\in\RR^{d_A}$ and $\mathrm{IA}_B\in\RR^{d_B}$:
\begin{enumerate}[noitemsep,label=\arabic*.]
\item \textbf{Pre-alignment:} share universal codebook $\mathcal{C}=\{g_1,\ldots,g_k\}$
  (VSA quasi-orthogonal vectors). Project to local spaces: $A\in\RR^{d_A\times k}$,
  $B\in\RR^{d_B\times k}$.
\item \textbf{Procrustes:} $M^*=\arg\min_{M^\top M=I}\norm{MA-B}_F$.
  Solution: $\mathrm{SVD}(BA^\top)=U\Sigma V^\top$, $M=UV^\top$.
\item \textbf{Transform:} $T_B=MT_AM^\dagger$.
\item \textbf{Apply:} $s_B'=T_B(s_B)$.
\end{enumerate}
For LoRA: $T_A=U_AV_A^\top\Rightarrow T_B.U=MU_A$, $T_B.V=V_AM^\dagger$ ---
rank $r$ is preserved.

%% ============================================================
\section{Engineering}
%% ============================================================

\paragraph{Binary format (.polydim V5, LoRA standard).}
$T=W_0+UV^\top$, $U,V\in\RR^{N\times r}$.
Dense (V2): $400\,\mathrm{MB}/\mathrm{transform}$ ($N=10{,}000$).
LoRA (V5): $\approx2.56\,\mathrm{MB}/\mathrm{transform}$ at $r=64$
--- a $156\times$ reduction.

\paragraph{Sparse execution.}
$A(\mathrm{DIM}_i)<\tau\Rightarrow$ skip tensor block.
Average complexity: $O(k_{\mathrm{active}}\cdot N\cdot r)$,
$k_{\mathrm{active}}\ll 9$.

%% ============================================================
\section{Human Interface: Flutter Isomorphism}
%% ============================================================

Flutter's widget tree admits a direct algebraic isomorphism to
\textsc{Polydim}'s composition algebra. The functor
$\varphi\colon\mathcal{T}\to\mathcal{W}$:
\begin{align*}
\varphi(\Compose(T_1,T_2))        &\cong \texttt{Column([$\varphi(T_1),\varphi(T_2)$])}\\
\varphi(\Mix(\alpha,T_1,\beta,T_2))&\cong \texttt{Stack([Opacity($\alpha$,$\varphi(T_1)$),$\ldots$])}\\
\varphi(\Fixpoint(T,\varepsilon))  &\cong \texttt{StatefulWidget(tol:$\varepsilon$)}
\end{align*}
When a user interacts with $\Project(\mathit{state},\mathsf{DIM\_FLUTTER})$,
the action generates $\delta_v\in\RR^N$ injected directly into the geometric
state: 2~steps vs.\ 8 in current AI-UI pipelines.

%% ============================================================
\section{Related Work}
%% ============================================================

\paragraph{Tensor frameworks.}
PyTorch~\citep{paszke2019pytorch}, JAX~\citep{bradbury2018jax}, and
TensorFlow~\citep{abadi2016tensorflow} use tensors as \emph{data} flowing
through sequential host-language programs. \textsc{Polydim} uses
transformations \emph{as} the program.

\paragraph{Agent communication.}
MCP~\citep{anthropic2024mcp} and A2A~\citep{google2025a2a} rely on JSON-RPC,
perpetuating the impedance mismatch. FIPA~\citep{fipa2002acl} grounds
communication in speech act theory. \textsc{Polydim} eliminates all three
layers via $T_B=MT_AM^\dagger$.

\paragraph{Categorical programming languages.}
Haskell~\citep{jones2003haskell}, Idris~\citep{brady2013idris},
and Agda~\citep{norell2007agda} apply categorical semantics to type theory of
sequential programs. Categorical Deep Learning~\citep{gavranovic2024categorical}
formalizes neural architectures as functors between parameter and data
categories; Backprop as Functor~\citep{fong2019backprop} characterizes supervised
learning compositionally. \textsc{Polydim} extends both to the full language-design
level, with $\Project$ as the canonical functor. Concurrently,
\citet{suarez2026invariant} provide an architectural proof that semantic features
occupy invariant linear subspaces, directly grounding~H2 of Theorem~\ref{thm:functor},
and \citet{oneill2025endofunctor} formalize self-attention as a parametric
endofunctor in $\mathrm{Para}(\mathrm{Vect})$, validating the categorical
framework of \textsc{Polydim}.

\paragraph{Vector Symbolic Architectures.}
VSA research~\citep{plate1995,kanerva2009,frady2021} establishes algebraic
operations preserving semantic structure. Quasi-orthogonality founds $\Mix$
as sound if/else replacement; VSA binding underlies $\GeoID$-based
cross-model identity.

\paragraph{Compiler infrastructure.}
MLIR~\citep{lattner2021mlir} provides multi-level IR. \textsc{Polydim} is
a \emph{Categorical IR}: programs are morphisms in $\mathcal{G}$;
compilation ($\Project$) is a functor.

\paragraph{Representation engineering.}
\citet{turner2023activation} and \citet{zou2023representation} demonstrate
activation-space interventions. \textsc{Polydim} generalizes both: $A\colon
D\to[0,1]$ exposes activation weights as first-class program parameters.

%% ============================================================
\section{Discussion}
%% ============================================================

\paragraph{What \textsc{Polydim} is not.}
Not a tensor framework (PyTorch programs compute \emph{with} tensors;
\textsc{Polydim} programs \emph{are} tensors). Not a neural architecture.
Not a new attention mechanism. Not a prompt language.

\paragraph{The assembly language parallel.}
Assembly formalized binary arithmetic programming without inventing the
transistor. \textsc{Polydim} formalizes geometric computation programming
without inventing the transformer. Unlike Architecture-specific Assembly,
\textsc{Polydim} is model-agnostic via $\Align$ --- closer to LLVM~IR.

\paragraph{Open problems.}
\textbf{Q1} (Subspace emergence): Do the 9 native subspaces emerge from
k-means clustering of LLM activations?~\citep{suarez2026invariant} predict
yes for high-frequency features.
\textbf{Q2} (Distributed consistency): Geometric-CRDTs where mutations
converge via geodesic barycenter on $S^{N-1}$.
\textbf{Q3} (Turing completeness): Conjectured via $\Fixpoint$ over
contractive $T$.
\textbf{Q4} (Self-hosting): Can a \textsc{Polydim} program generate
\textsc{Polydim} programs?

%% ============================================================
\section{Conclusion}
%% ============================================================

We have presented \textsc{Polydim}, a programming language whose fundamental
unit is a geometric transformation $T\colon\RR^N\to\RR^N$. Four algebraic
theorems establish the formal foundation.
Theorem~\ref{thm:functor}---$\Project$ is a functor---is the central
mathematical contribution, grounded in the Invariant Subspace Necessity
theorem~\citep{suarez2026invariant} and the categorical endofunctor
framework~\citep{oneill2025endofunctor}. Hypothesis~H2, previously the
main open question, is now a consequence of the transformer architecture
itself. Corollaries~\ref{cor:portability} and~\ref{cor:pullback} guarantee
cross-executor portability and intersection types as categorical pullbacks ---
a live-bound type with no analogue in single-paradigm languages. The Python
bootstrap validates the framework with $29/29$ tests, $0.9993$ alignment score,
and $+64.5\%$ semantic improvement. The Rust core, LoRA binary format,
Procrustes-based ALIGN, and Flutter executor complete the engineering foundation.
The assembly language of the transformer era is here.

%% ============================================================
\bibliographystyle{plainnat}
\bibliography{polydim}
%% ============================================================
%% Add to polydim.bib before compiling:
%%   @article{suarez2026invariant,...} (see polydim_bib_addendum.bib)
%%   @article{oneill2025endofunctor,...} (see polydim_bib_addendum.bib)
%% ============================================================

%% ============================================================
\appendix
\section{Bootstrap V0.3 Metrics}
%% ============================================================

\begin{table}[h]
\centering
\small
\begin{tabular}{ll}
\toprule
Metric & Value \\
\midrule
Tests passing            & $29/29$ \\
align\_score (MODO\_H)   & $0.9993$ \\
Semantic sim.\ (MiniLM)  & $0.830$ \\
Semantic sim.\ (determ.) & $0.506$ \\
Semantic improvement     & $+64.5\%$ \\
\bottomrule
\end{tabular}
\caption{Bootstrap V0.3 verified metrics.}
\end{table}

\noindent Known limitations of the bootstrap (not of \textsc{Polydim}):
subspaces are string aliases (not geometric regions);
$\Compose$ and $\Mix$ not yet implemented as real matrix ops (resolved in
\texttt{polydim\_primitives\_v1.py});
$\Align$: Procrustes implemented in \texttt{polydim\_align\_v2.py};
Rust core: \texttt{polydim\_core.rs};
$\Project(\mathsf{DIM\_FLUTTER})$: \texttt{polydim\_flutter\_executor.dart}.

\end{document}
